% !TeX root = ../main.tex
\chapter{Giới thiệu}
\label{chap:introduction}

\section{Cây quyết định}

Cây quyết định là mô hình học có giám sát biểu diễn quá trình dự đoán bằng một
chuỗi điều kiện dạng ``nếu--thì''. Tại mỗi node, thuật toán chọn một feature và
một ngưỡng chia dữ liệu thành các nhóm có độ hỗn tạp nhỏ hơn. Với bài toán phân
loại, node lá trả về lớp chiếm ưu thế trong các mẫu huấn luyện đi đến lá đó.
Mô hình cây phân loại và cơ chế chọn phép chia được trình bày trong CART và tài
liệu scikit-learn~\cite{breiman-cart,sklearn-tree}.

Trong thí nghiệm này, mô hình sử dụng chỉ số Gini mặc định. Với tập mẫu tại
node $t$, Gini impurity được xác định bởi
\begin{equation}
  G(t)=1-\sum_{k=1}^{K}p(k\mid t)^2,
\end{equation}
trong đó $p(k\mid t)$ là tỷ lệ mẫu lớp $k$ tại node. Một phép chia tốt làm giảm
độ hỗn tạp có trọng số của hai node con.

Ưu điểm quan trọng của Decision Tree là không cần chuẩn hóa dữ liệu, có thể học
ranh giới phi tuyến và cho phép đọc trực tiếp các luật quyết định. Nhược điểm là
một cây không bị ràng buộc dễ tạo các nhánh rất đặc thù cho tập training, dẫn
đến phương sai cao và overfitting.

\section{Mục tiêu dự án}

Dự án áp dụng Decision Tree cho bài toán phân loại chẩn đoán khối u vú. Các
mục tiêu cụ thể gồm:

\begin{itemize}
  \item lựa chọn và mô tả một dataset công khai, phù hợp cho phân loại;
  \item xây dựng baseline trên một train/test split tái lập;
  \item đánh giá bằng accuracy, error rate, precision, recall, F1 và confusion
  matrix;
  \item phân tích cấu trúc, feature và các luật chia quan trọng của cây;
  \item thử độc lập ba cách cải thiện: giới hạn \texttt{max\_depth}, điều chỉnh
  \texttt{min\_samples\_leaf} và cost-complexity pruning;
  \item so sánh hiệu năng với độ phức tạp và chọn mô hình phù hợp nhất.
\end{itemize}

\section{Phạm vi và nguyên tắc đánh giá}

Mục tiêu không chỉ là tìm con số accuracy lớn nhất. Báo cáo xem xét error rate,
khả năng phát hiện lớp malignant, khoảng cách train--test và độ phức tạp của
cây. Hyperparameter được chọn bằng cross-validation trên tập training; tập test
chỉ dùng trong so sánh cuối cùng.
